% concatenated from Discriminant_pushed__kopia_.tex

\documentclass[11pt]{amsart}

\usepackage{amsmath,amssymb,amsthm,mathrsfs,enumitem}
\usepackage[colorlinks=true,linkcolor=blue,citecolor=blue,urlcolor=blue]{hyperref}

\numberwithin{equation}{section}
\emergencystretch=2em

\newtheorem{theorem}{Theorem}
\newtheorem{proposition}[theorem]{Proposition}
\newtheorem{lemma}[theorem]{Lemma}
\newtheorem{corollary}[theorem]{Corollary}
\newtheorem{remark}[theorem]{Remark}
\newtheorem{definition}[theorem]{Definition}
\newtheorem{conjecture}[theorem]{Conjecture}
\newtheorem{problem}[theorem]{Problem}
\newtheorem{example}[theorem]{Example}

\newcommand{\bC}{\mathbb C}
\newcommand{\bR}{\mathbb R}
\newcommand{\D}{\mathcal D}
\newcommand{\Res}{\operatorname{Res}}
\newcommand{\disc}{\operatorname{disc}}
\newcommand{\Sym}{\operatorname{Sym}}
\newcommand{\barx}{\overline{x}}

\title[Discriminants and square-graph cones]
{Discriminants of derivatives and the square-graph cone}

%\author{Vladimir Petrov Kostov}
%\address{Universit\'e C\^ote d'Azur, Laboratoire J.A.~Dieudonn\'e,
%Parc Valrose, 06108 Nice Cedex 2, France}
%\email{kostov@math.unice.fr}

\author{Boris Shapiro}
\address{Department of Mathematics, Stockholm University,
S--10691 Stockholm, Sweden}
\email{shapiro@math.su.se}

\begin{document}

\begin{abstract}
Let $P$ be a monic polynomial of degree $n$ with roots $x_1,\ldots,x_n$.
We study the discriminants of the derivatives $P^{(k)}$ as symmetric
translation-invariant polynomials in the original roots.  Alexandersson and
Shapiro conjectured that every such discriminant belongs to the cone generated
by symmetrized graph monomials with even edge multiplicities.  We obtain a
sharp positive/negative picture in the first terminal cases.  For the terminal
cubic family $k=n-3$ we prove the conjecture for every $n\ge3$, and for
$n\ge5$ obtain the three-graph formula conjectured in \cite[Example~3]{AS}.
For the terminal quartic family we show, by exact finite computation, that
$\disc(P^{(n-4)})$ belongs to the square-graph cone if and only if
$4\le n\le22$.  Thus the general square-graph cone conjecture fails already
at $(n,k)=(23,19)$ and fails for every $k=n-4$ with $n\ge23$.  The negative
result is certified by an explicit linear functional which is nonnegative on
every degree-$12$ square-graph generator and strictly negative on the quartic
discriminant.  We also record central-moment formulas, the subset-average and
finite Appell structure of normalized terminal polynomials, and the explicit
quintic member.
\end{abstract}

\subjclass[2020]{Primary 26C10, 05C31; Secondary 12D10, 05E05, 14M25}

\keywords{discriminant, derivative of a polynomial, symmetrized graph monomial,
sum of squares, square-graph cone, central moments}

\maketitle

\section{Introduction}

Let
\[
P(t)=\prod_{i=1}^n(t-x_i)=t^n+a_1t^{n-1}+\cdots+a_n
\]
be a monic polynomial.  For $0\le k\le n-2$ put
\[
\D_k^{(n)}:=\disc(P^{(k)}),
\]
where the discriminant is taken with the usual leading-coefficient
normalization.
Equivalently,
\[
\D_k^{(n)}=(-1)^{m(m-1)/2}c^{-1}\Res(P^{(k)},P^{(k+1)}),\qquad
m=n-k,
\]
where $c$ is the leading coefficient of $P^{(k)}$.

The polynomial $\D_k^{(n)}$ is symmetric and translation-invariant in
$x_1,\ldots,x_n$.  Hence it lies in the algebra generated by the pairwise
root differences $x_i-x_j$.  It is natural to ask whether it admits a
representation by manifestly nonnegative expressions in these differences.

The following strengthening of a question of Sottile and Mukhin was formulated
in \cite{AS}.

\begin{conjecture}[square-graph cone conjecture]\label{conj:squaregraph}
For every $n\ge2$ and every $0\le k\le n-2$, the polynomial $\D_k^{(n)}$
belongs to the convex cone generated by symmetrized graph monomials whose
edges all have even multiplicity.
Equivalently, $\D_k^{(n)}$ is a positive linear combination of symmetric sums
of products
\[
\prod_{(i,j)\in E(G)}(x_i-x_j)^2 .
\]
\end{conjecture}

For $k=0$ this is the usual discriminant formula.  For the first derivative,
the weaker assertion that $\disc(P')$ is a sum of squares in the differences of
the roots was proved by Sanyal, Sturmfels and Vinzant via the entropic
discriminant \cite[Corollary~14]{Sanyal}.  Alexandersson and Shapiro found
considerable evidence for Conjecture~\ref{conj:squaregraph}: in particular,
\cite[Example~3]{AS} proposed an explicit three-square-graph formula for the
terminal cubic family, and their Examples~4--5 give positive terminal quartic
certificates for $n=5,6$.

The main point of the present note is that this positive pattern has a sharp
boundary.  We work in the terminal range, where the order
\[
        r=n-k
\]
of the remaining derivative is fixed.  The main contributions are the
following.
\begin{enumerate}[label={\rm(\alph*)}]
\item We give closed central-moment formulas for the terminal quadratic,
cubic and quartic discriminants.
\item We prove a positive square-graph expansion for the terminal cubic
family.  Consequently Conjecture~\ref{conj:squaregraph} holds for
$\disc(P^{(n-3)})$ for every $n\ge3$.
\item For $n\ge5$ we compress the cubic certificate to three graph types.
After accounting for symmetrization conventions, the resulting coefficients
are exactly those conjectured in \cite[Example~3]{AS}; thus the earlier finite
verification becomes an all-parameter proof.
\item For the terminal quartic family we obtain a sharp classification:
$\disc(J_{n,4})$ belongs to the degree-$12$ square-graph cone exactly for
$4\le n\le22$.  For $n\ge23$ we give an explicit separating linear functional
which is nonnegative on every square-graph generator and strictly negative on
$\disc(J_{n,4})$.  This disproves Conjecture~\ref{conj:squaregraph}, with the
first terminal-quartic counterexample at $(n,k)=(23,19)$.
\item We introduce normalized terminal polynomials $J_{n,r}$ for arbitrary
fixed terminal order $r$.  They are uniform averages of the characteristic
polynomials of $r$-subsets of the centered roots and satisfy the finite Appell
relation $J'_{n,r}=rJ_{n,r-1}$.  We also write down the quintic member
explicitly.
\item We show that fixed degree reduces the square-graph cone to finitely many
reduced graph types.  The quartic theorem shows that this finiteness should be
used on the \emph{dual} side as well: stable separating functionals, rather
than uniformly positive primal certificates, become the natural objects.
\end{enumerate}
Thus the cubic family is uniformly positive, whereas the very next terminal
order crosses out of the square-graph cone at the finite threshold $n=23$.

\section{Graph monomials and the ambient vector space}

We recall the basic language.  If $G$ is a loopless multigraph with vertex set
$\{1,\ldots,n\}$, define
\[
M_G(x_1,\ldots,x_n)=\prod_{(i,j)\in E(G)}(x_i-x_j),
\]
where an orientation may be chosen for each edge.  Its symmetrization is
\[
\widetilde G:=\sum_{\sigma\in S_n}M_G(x_{\sigma(1)},\ldots,x_{\sigma(n)}).
\]
If every edge of $G$ occurs with even multiplicity, then $\widetilde G$ is a
sum of squares, since it is a sum of squares of monomials in differences.
Such graphs will be called \emph{square graphs}.

Let $\mathrm{PST}_{n,d}$ denote the vector space of homogeneous symmetric
translation-invariant polynomials of degree $d$ in $n$ variables.  It is
classical that this space is spanned by symmetrized graph monomials with $d$
edges; see \cite{AS} and the references therein.  Since
\[
\deg \D_k^{(n)}=(n-k)(n-k-1),
\]
Conjecture~\ref{conj:squaregraph} says that a distinguished element of
$\mathrm{PST}_{n,(n-k)(n-k-1)}$ lies in the square-graph cone.

\section{The quadratic terminal case}

Let
\[
\barx=\frac1n\sum_{i=1}^n x_i,
\qquad
z_i=x_i-\barx,
\qquad
p_r=\sum_{i=1}^n z_i^r .
\]
Thus $p_1=0$.

\begin{proposition}\label{prop:quadratic}
For every $n\ge2$,
\[
\D_{n-2}^{(n)}
=(n-1)((n-2)!)^2\sum_{1\le i<j\le n}(x_i-x_j)^2.
\]
Equivalently,
\[
\D_{n-2}^{(n)}=n(n-1)((n-2)!)^2p_2.
\]
\end{proposition}

\begin{proof}
The polynomial $P^{(n-2)}$ is quadratic.  Its coefficients are
\[
P^{(n-2)}(t)=\frac{n!}{2}t^2+(n-1)!a_1t+(n-2)!a_2.
\]
Therefore
\[
\disc(P^{(n-2)})=((n-1)!a_1)^2-2n!(n-2)!a_2.
\]
Using $a_1=-\sum_i x_i$ and $a_2=\sum_{i<j}x_ix_j$, this becomes
\[
(n-1)((n-2)!)^2\left((n-1)\Big(\sum_i x_i\Big)^2
-2n\sum_{i<j}x_ix_j\right).
\]
The expression in parentheses is
$\sum_{i<j}(x_i-x_j)^2=np_2$, which proves the claim.
\end{proof}

This proves Conjecture~\ref{conj:squaregraph} in the terminal quadratic case.
It is also the formula appearing as Example~2 in \cite{AS}, with the same
normalization of the discriminant.

\section{The cubic terminal case}

The next terminal case is already more interesting.  It gives a compact formula
in terms of the second and third central moments of the roots.

\begin{theorem}\label{thm:cubic}
For every $n\ge3$,
\[
\boxed{
\D_{n-3}^{(n)}=\frac{n!((n-2)!)^3}{12}
\left(p_2^3-\frac{n(n-1)}{(n-2)^2}p_3^2\right).}
\]
In particular, if the roots $x_1,\ldots,x_n$ are real, then
\[
 p_2^3\ge \frac{n(n-1)}{(n-2)^2}p_3^2 .
\]
\end{theorem}

\begin{proof}
Translate the variable so that the roots become $z_i=x_i-\barx$ and
$\sum_i z_i=0$.  Translation does not change the discriminant.
The polynomial $P$ then has the form
\[
P(t)=t^n+e_2t^{n-2}-e_3t^{n-3}+\cdots,
\]
where, by Newton's identities,
\[
e_2=-\frac{p_2}{2},\qquad e_3=\frac{p_3}{3}.
\]
After differentiating $n-3$ times, we obtain the depressed cubic
\[
P^{(n-3)}(t)=\frac{n!}{6}t^3+(n-2)!e_2t-(n-3)!e_3.
\]
For a depressed cubic $At^3+Bt+C$, the discriminant is
\[
-4AB^3-27A^2C^2.
\]
Substituting
\[
A=\frac{n!}{6},\qquad
B=-\frac{(n-2)!}{2}p_2,
\qquad
C=-\frac{(n-3)!}{3}p_3
\]
gives
\[
\D_{n-3}^{(n)}
=\frac{n!((n-2)!)^3}{12}p_2^3
-\frac{(n!)^2((n-3)!)^2}{12}p_3^2.
\]
Since
\[
\frac{(n!)^2((n-3)!)^2}{n!((n-2)!)^3}
=\frac{n(n-1)}{(n-2)^2},
\]
the stated formula follows.  If all $x_i$ are real, then $P^{(n-3)}$ is a real
cubic with three real roots by Rolle's theorem, and hence its discriminant is
nonnegative.  This gives the displayed moment inequality.
\end{proof}

\begin{remark}
The inequality in Theorem~\ref{thm:cubic} is the sharp finite version of the
usual bound on skewness for a distribution supported on $n$ equally weighted
points.  The normalization here is especially natural because it is forced by
the discriminant of $P^{(n-3)}$.
\end{remark}

\begin{proposition}\label{prop:equality}
Assume $x_1,\ldots,x_n\in\bR$ are not all equal.  Equality in
\[
 p_2^3\ge \frac{n(n-1)}{(n-2)^2}p_3^2
\]
holds if and only if, after a permutation, the centered roots have the form
\[
z_1=\cdots=z_{n-1}=a,
\qquad
z_n=-(n-1)a
\]
for some nonzero real number $a$, or the same configuration with $a$ replaced
by $-a$.
\end{proposition}

\begin{proof}
Equality is equivalent to $\D_{n-3}^{(n)}=0$, i.e. to $P^{(n-3)}$ having a
multiple root.  Since $P$ has real roots, $P^{(n-3)}$ has three real roots
counted with multiplicity.  A depressed cubic with a multiple root has root
multiset $\{u,u,-2u\}$ for some $u\in\bR$.

Equivalently, the extremal problem for $p_3$ under the constraints
$\sum_i z_i=0$ and $\sum_i z_i^2=p_2$ has a Lagrange multiplier equation
\[
3z_i^2=\lambda+2\mu z_i.
\]
Thus all $z_i$ assume at most two distinct values.  If these values have
multiplicities $r$ and $n-r$, the constraints give
\[
|p_3|=\frac{|n-2r|}{\sqrt{r(n-r)n}}p_2^{3/2}.
\]
This is maximal precisely for $r=1$ or $r=n-1$, giving the stated equality
case.
\end{proof}

\section{A quartic terminal formula}

The next case is still computable in closed form.  It is useful because it
is the first terminal case involving the fourth central moment.

\begin{theorem}\label{thm:quartic}
Let $n\ge4$.  With the notation
\[
p_r=\sum_{i=1}^n (x_i-\bar x)^r ,
\qquad
\bar x=\frac1n\sum_i x_i ,
\]
one has
\[
\D_{n-4}^{(n)}
=\disc\left(
A t^4+B t^2+C t+D
\right),
\]
where
\[
\begin{aligned}
A&=\frac{n!}{24},
&B&=-\frac{(n-2)!}{4}p_2,\\
C&=-\frac{(n-3)!}{3}p_3,
&D&=(n-4)!\left(\frac{p_2^2}{8}-\frac{p_4}{4}\right).
\end{aligned}
\]
Equivalently,
\[
\begin{aligned}
\D_{n-4}^{(n)}
={}&256A^3D^3-128A^2B^2D^2+144A^2BC^2D-27A^2C^4\\
&\quad +16AB^4D-4AB^3C^2 .
\end{aligned}
\]
In particular, for real roots $x_1,\ldots,x_n$ this gives the sharp
finite moment inequality
\[
256A^3D^3-128A^2B^2D^2+144A^2BC^2D-27A^2C^4
+16AB^4D-4AB^3C^2\ge0 .
\]
\end{theorem}

\begin{proof}
After translating the roots we may assume $\sum_i x_i=0$.  Newton's identities
give
\[
e_2=-\frac{p_2}{2},\qquad
e_3=\frac{p_3}{3},\qquad
e_4=\frac{p_2^2}{8}-\frac{p_4}{4}.
\]
Hence
\[
P(t)=t^n+e_2t^{n-2}-e_3t^{n-3}+e_4t^{n-4}+\cdots .
\]
Differentiating $n-4$ times gives
\[
P^{(n-4)}(t)
=
\frac{n!}{24}t^4+\frac{(n-2)!}{2}e_2t^2
-(n-3)!e_3t+(n-4)!e_4,
\]
which is precisely the stated quartic.  The discriminant formula for
$At^4+Bt^2+Ct+D$ is
\[
256A^3D^3-128A^2B^2D^2+144A^2BC^2D-27A^2C^4
+16AB^4D-4AB^3C^2 .
\]
Finally, if the $x_i$ are real, then all roots of $P^{(n-4)}$ are real by
Rolle's theorem; therefore its discriminant is nonnegative.
\end{proof}

\begin{corollary}[normalized quartic form]\label{cor:normalized-quartic}
Let
\[
q_{n,4}(t)=t^4+\alpha t^2+\beta t+\gamma,
\]
where
\[
\begin{aligned}
\alpha&=-\frac{6p_2}{n(n-1)},
&\beta&=-\frac{8p_3}{n(n-1)(n-2)},\\
\gamma&=\frac{3(p_2^2-2p_4)}{n(n-1)(n-2)(n-3)}.
\end{aligned}
\]
Then, after centering the roots,
\[
P^{(n-4)}(t)=\frac{n!}{24}\,q_{n,4}(t)
\]
and hence
\[
\D_{n-4}^{(n)}=\left(\frac{n!}{24}\right)^6\disc(q_{n,4}).
\]
In particular, for real roots the depressed quartic $q_{n,4}$ is hyperbolic
and
\[
\begin{aligned}
0\le \disc(q_{n,4})={}&256\gamma^3-128\alpha^2\gamma^2
+144\alpha\beta^2\gamma-27\beta^4\\
&\hspace{1.3cm}+16\alpha^4\gamma-4\alpha^3\beta^2 .
\end{aligned}
\]
If $p_3=0$, this specializes to the factorization
\[
\disc(q_{n,4})=16\gamma(\alpha^2-4\gamma)^2 .
\]
\end{corollary}

\begin{proof}
Divide the quartic in Theorem~\ref{thm:quartic} by its leading coefficient
$n!/24$.  The formula for the discriminant follows from the homogeneity rule
$\disc(cq)=c^{2r-2}\disc(q)$ for a degree $r$ polynomial.  The final
factorization is the standard discriminant factorization for a biquadratic
quartic $t^4+\alpha t^2+\gamma$.
\end{proof}

\begin{remark}
The quartic formula gives a natural continuation of the cubic skewness
inequality.  In the language of finite probability distributions with equal
weights it is an algebraic inequality among the variance, skewness, and fourth
central moment.  Unlike the cubic case, the equality locus is no longer only a
two-point configuration; it corresponds to configurations for which
$P^{(n-4)}$ has a multiple root.  The next section shows that, despite this
nonnegativity, the normalized discriminant eventually leaves the much smaller
square-graph cone.
\end{remark}

\section{A sharp quartic obstruction}\label{sec:quartic-obstruction}

Put
\[
   \Delta_{n,4}:=\disc(J_{n,4})
   =\left(\frac{24}{n!}\right)^6\D_{n-4}^{(n)},
\]
and let $\mathcal C_{n,12}\subset\mathrm{PST}_{n,12}$ denote the cone
generated by degree-$12$ square graphs.  Since the scalar relating
$\Delta_{n,4}$ and $\D_{n-4}^{(n)}$ is positive, the two polynomials have the
same cone-membership behavior.

For a homogeneous symmetric form $F$ of degree $12$ write
\[
       F=\sum_{\lambda\vdash12}c_\lambda(F)m_\lambda,
\]
where $m_\lambda$ is the monomial symmetric polynomial (terms with
$\ell(\lambda)>n$ are absent).  For $n\ge12$ define the linear functional
\begin{align}\label{eq:quartic-separator}
\Lambda_n(F)={}&96\bigl(
 c_{(12)}+c_{(10,2)}-c_{(9,3)}+c_{(8,4)}+c_{(8,2,2)}
 -c_{(7,5)}-c_{(7,3,2)} \notag\\
&\hspace{31mm}+c_{(6,6)}+c_{(5,4,3)}\bigr)
 +32c_{(4,4,4)}+(n-3)c_{(3,3,3,3)}.
\end{align}

\begin{theorem}[sharp terminal-quartic threshold]\label{thm:quartic-threshold}
For every integer $n\ge4$,
\[
             \Delta_{n,4}\in\mathcal C_{n,12}
             \quad\Longleftrightarrow\quad 4\le n\le22.
\]
Consequently Conjecture~\ref{conj:squaregraph} is false.  More precisely,
\[
       \D_{n-4}^{(n)}\notin\mathcal C_{n,12}
       \qquad(n\ge23),
\]
so $(n,k)=(23,19)$ is the first counterexample inside the terminal quartic
family.
\end{theorem}

The proof is computer-assisted but entirely exact.  We give the algebraic
reduction and the separating functional explicitly; the finite checks are
contained in the ancillary verifier described below.

Let $H$ be a loopless multigraph with six edges counted with multiplicity,
and let $2H$ denote the square graph obtained by doubling every edge.  If
$H$ has $v$ nonisolated labelled vertices, define its injective orbit sum by
\[
 \mathcal O_H^{(n)}=
 \sum_{\phi:[v]\hookrightarrow[n]}
 \prod_{uv\in E(H)}(x_{\phi(u)}-x_{\phi(v)})^{2m_{uv}}.
\]
After adjoining $n-v$ isolated vertices, the full $S_n$-symmetrization is
$(n-v)!\mathcal O_H^{(n)}$; hence the two conventions generate the same ray.
Every degree-$12$ square graph arises in this way.

For completeness, there is a particularly simple exact formula for the
monomial-symmetric coordinates of these orbit sums.  Put
\[
 Q_H(y_1,\ldots,y_v)=
 \prod_{uv\in E(H)}(y_u-y_v)^{2m_{uv}}
 =\sum_{\alpha\in\mathbb Z_{\ge0}^v}q_\alpha y^\alpha.
\]
For $\lambda\vdash12$, let $\alpha^+$ denote the positive entries of $\alpha$
and set
\[
 b_\lambda(H)=
 \left(\prod_{a\ge1}\operatorname{mult}_a(\lambda)!\right)
 \sum_{\operatorname{sort}(\alpha^+)=\lambda}q_\alpha .
\]
Then, whenever $n\ge v$,
\begin{equation}\label{eq:orbit-monomial-coordinate}
 [m_\lambda]\mathcal O_H^{(n)}
 =(n-\ell(\lambda))_{\downarrow\,v-\ell(\lambda)}b_\lambda(H).
\end{equation}
Thus every coordinate is an explicit integer polynomial in $n$.

\begin{lemma}[dual certificate]\label{lem:quartic-dual}
For every degree-$12$ square graph and every $n\ge12$ one has
\[
                   \Lambda_n(\widetilde G)\ge0.
\]
On the other hand,
\begin{equation}\label{eq:quartic-target-separation}
 \Lambda_n(\Delta_{n,4})=
 -\frac{2654208\,Q(n)}
 {n^6(n-3)^2(n-2)^4(n-1)^5},
\end{equation}
where
\[
 Q(n)=n^6-36n^5+376n^4-1920n^3+5635n^2-9456n+6912.
\]
In particular, $\Lambda_n(\Delta_{n,4})<0$ for every $n\ge23$.
\end{lemma}

\begin{proof}
Dividing all even edge multiplicities by two identifies the degree-$12$
square-graph types with loopless multigraphs having six edges.  There are
$212$ such reduced isomorphism types; this is the enumeration already used in
\cite{AS}.  The ancillary exact verifier independently reconstructs the same
list.  For each type it expands \eqref{eq:orbit-monomial-coordinate}, applies
\eqref{eq:quartic-separator}, and substitutes $n=m+12$.  In all $212$ cases
the resulting polynomial in $m$ has nonnegative integer coefficients (in
$53$ cases it is identically zero).  Hence $\Lambda_n$ is nonnegative on every
generator for all $n\ge12$.

Independently expanding the discriminant of the normalized quartic in
Corollary~\ref{cor:normalized-quartic} in the monomial-symmetric basis gives
\eqref{eq:quartic-target-separation}.  Finally,
\[
\begin{aligned}
Q(m+23)={}&m^6+102m^5+4171m^4+85572m^3\\
          &\quad+884074m^2+3748560m+957456,
\end{aligned}
\]
whose coefficients are all positive.  Therefore $Q(n)>0$ for $n\ge23$, and
the last assertion follows.
\end{proof}

\begin{proof}[Proof of Theorem~\ref{thm:quartic-threshold}]
The non-membership for $n\ge23$ follows immediately from
Lemma~\ref{lem:quartic-dual} by separation.

For $4\le n\le22$ we use exact positive certificates in the injective orbit
normalization above.  They are stored in the ancillary file
\texttt{quartic\_positive\_certificates.json}; every coefficient is a positive
rational number.  The exact verifier \texttt{verify\_quartic\_separator.py}
checks coefficientwise equality in the monomial-symmetric basis, discarding
partitions of length $>n$ and allowing only graph types with at most $n$
nonisolated vertices.  No floating-point arithmetic enters this verification.
The numbers of graph rays used are
\[
\begin{array}{c|ccccccccc}
 n&4&5&6&7&8&9&10&11&12\text{--}22\\ \hline
 \#&1&3&6&8&10&11&11&12&13.
\end{array}
\]
Thus $\Delta_{n,4}\in\mathcal C_{n,12}$ throughout $4\le n\le22$, completing
the proof.
\end{proof}

\begin{remark}
The obstruction begins at $n=23$, just before the convenient stable range
$n\ge24$ associated with degree $12$.  Thus the failure is not an artifact of
passing to an eventually stable ambient basis.  The functional
\eqref{eq:quartic-separator} already separates the whole infinite tail.
\end{remark}

\section{Universal terminal polynomials}

For $m\ge j$ write
\[
(m)_{\downarrow j}=m(m-1)\cdots(m-j+1)
\]
for the falling factorial.  The preceding examples are instances of the
following normalized terminal principle.

\begin{theorem}[subset-average and Appell structure]\label{thm:terminal-normalized}
Fix $r\ge2$ and $n\ge r$.  After translating the roots so that
$\sum_i x_i=0$, put
\[
J_{n,r}(t)=\frac{r!}{n!}P^{(n-r)}(t).
\]
Then one has the subset-average representation
\begin{equation}\label{eq:subset-average}
J_{n,r}(t)=\binom nr^{-1}
\sum_{\substack{S\subseteq\{1,\ldots,n\}\\ |S|=r}}
\prod_{i\in S}(t-x_i).
\end{equation}
Equivalently,
\begin{equation}\label{eq:terminal-coefficients}
J_{n,r}(t)
=t^r+\sum_{j=2}^r(-1)^j
\frac{\binom rj}{\binom nj}e_jt^{r-j}
=t^r+\sum_{j=2}^r(-1)^j
\frac{(r)_{\downarrow j}}{(n)_{\downarrow j}}e_jt^{r-j}.
\end{equation}
Consequently $J_{n,r}$ is a universal polynomial in
$p_2,p_3,\ldots,p_r$, and
\[
\D_{n-r}^{(n)}=\left(\frac{n!}{r!}\right)^{2r-2}
\disc(J_{n,r}).
\]
Moreover, with $J_{n,1}(t)=t$ and $J_{n,0}(t)=1$, the terminal polynomials
form a finite Appell chain:
\begin{equation}\label{eq:appell}
J'_{n,r}(t)=rJ_{n,r-1}(t),\qquad 1\le r\le n.
\end{equation}
In particular,
\begin{equation}\label{eq:terminal-resultant}
\disc(J_{n,r})
=(-1)^{r(r-1)/2}r^r\Res(J_{n,r},J_{n,r-1}).
\end{equation}
If the roots $x_i$ are real, then all $J_{n,r}$ are real-rooted, and the
zeros of consecutive members interlace.
\end{theorem}

\begin{proof}
Differentiate the product $P(t)=\prod_{i=1}^n(t-x_i)$ exactly $n-r$
times.  For each $r$-subset $S$, the product
$\prod_{i\in S}(t-x_i)$ occurs $(n-r)!$ times, whence
\[
P^{(n-r)}(t)=(n-r)!\sum_{|S|=r}\prod_{i\in S}(t-x_i).
\]
Multiplication by $r!/n!=\bigl(\binom nr(n-r)!\bigr)^{-1}$ gives
\eqref{eq:subset-average}.  The coefficient of $(-1)^je_jt^{r-j}$ in
the average is the probability that a fixed $j$-subset is contained in a
uniformly chosen $r$-subset, namely
$\binom{n-j}{r-j}/\binom nr=\binom rj/\binom nj$; this proves
\eqref{eq:terminal-coefficients}.  Since $e_1=p_1=0$, Newton's identities
express every $e_j$ as a polynomial in $p_2,\ldots,p_j$.

The discriminant scaling formula gives the displayed normalization of
$\D_{n-r}^{(n)}$.  Differentiating the definition of $J_{n,r}$ yields
\[
J'_{n,r}=\frac{r!}{n!}P^{(n-r+1)}=rJ_{n,r-1},
\]
which is \eqref{eq:appell}.  Since $J_{n,r}$ is monic,
\[
\disc(J_{n,r})=(-1)^{r(r-1)/2}\Res(J_{n,r},J'_{n,r}),
\]
and \eqref{eq:terminal-resultant} follows from the homogeneity of the
resultant in its second argument.  Finally, repeated Rolle interlacing for
successive derivatives of a real-rooted polynomial gives real-rootedness and
interlacing of the entire chain.
\end{proof}

The first three normalized polynomials are
\[
J_{n,2}(t)=t^2-\frac{p_2}{n(n-1)},
\]
\[
J_{n,3}(t)=t^3-\frac{3p_2}{n(n-1)}t
-\frac{2p_3}{n(n-1)(n-2)},
\]
and $J_{n,4}=q_{n,4}$ from Corollary~\ref{cor:normalized-quartic}.  Notice in
particular that $q'_{n,4}=4J_{n,3}$: the cubic moment inequality of
Theorem~\ref{thm:cubic} is therefore also the discriminant condition for the
critical points of the terminal quartic.  The next case is still short enough
to be useful.

\begin{proposition}[the quintic terminal polynomial]\label{prop:quintic-terminal}
For $n\ge5$,
\[
\begin{aligned}
J_{n,5}(t)={}&t^5-\frac{10p_2}{n(n-1)}t^3
-\frac{20p_3}{n(n-1)(n-2)}t^2\\
&\quad+\frac{15p_2^2-30p_4}{n(n-1)(n-2)(n-3)}t\\
&\quad+\frac{20p_2p_3-24p_5}
{n(n-1)(n-2)(n-3)(n-4)} .
\end{aligned}
\]
Thus, for real roots, the discriminant of this quintic is a nonnegative
universal polynomial in $p_2,p_3,p_4,p_5$.
\end{proposition}

\begin{proof}
Apply Theorem~\ref{thm:terminal-normalized} with $r=5$ and use Newton's
identities in the centered case:
\[
e_2=-\frac{p_2}{2},\qquad e_3=\frac{p_3}{3},\qquad
e_4=\frac{p_2^2}{8}-\frac{p_4}{4},\qquad
e_5=\frac{p_5}{5}-\frac{p_2p_3}{6}.
\]
Substitution gives the displayed expression.
\end{proof}

\begin{corollary}[an odd quintic slice]\label{cor:quintic-odd-slice}
Assume $n\ge5$ and $p_3=p_5=0$.  Set
\[
 a=-\frac{10p_2}{n(n-1)},\qquad
 b=\frac{15(p_2^2-2p_4)}{n(n-1)(n-2)(n-3)}.
\]
Then
\[
 J_{n,5}(t)=t(t^4+at^2+b)
\]
and
\begin{equation}\label{eq:quintic-odd-discriminant}
       \disc(J_{n,5})=16b^3(a^2-4b)^2.
\end{equation}
Equivalently,
\[
\disc(J_{n,5})=
\frac{86400000\,(p_2^2-2p_4)^3
\bigl((n^2-11n+15)p_2^2+3n(n-1)p_4\bigr)^2}
{n^7(n-1)^7(n-2)^5(n-3)^5}.
\]
If the original roots are real, then
\begin{equation}\label{eq:odd-slice-moment-bounds}
 p_4\le\frac12p_2^2,
 \qquad
 (n^2-11n+15)p_2^2+3n(n-1)p_4\ge0.
\end{equation}
Both boundary factors in \eqref{eq:quintic-odd-discriminant} have the expected
multiple-root interpretation for the terminal quintic.
\end{corollary}

\begin{proof}
Under $p_3=p_5=0$, Proposition~\ref{prop:quintic-terminal} gives the displayed
odd quintic.  Since
\[
 \disc\bigl(t(t^4+at^2+b)\bigr)
 =\Res(t,t^4+at^2+b)^2\disc(t^4+at^2+b),
\]
and $\Res(t,t^4+at^2+b)=b$, Corollary~\ref{cor:normalized-quartic} gives
\eqref{eq:quintic-odd-discriminant}.  Moreover
\[
 a^2-4b=
 \frac{40\bigl((n^2-11n+15)p_2^2+3n(n-1)p_4\bigr)}
 {n^2(n-1)^2(n-2)(n-3)}.
\]
For real original roots the polynomial $J_{n,5}$ is real-rooted by
Theorem~\ref{thm:terminal-normalized}.  Hence its nonzero roots occur as
$\pm u,\pm v$, so $b=u^2v^2\ge0$ and
$a^2-4b=(u^2-v^2)^2\ge0$.  This proves
\eqref{eq:odd-slice-moment-bounds}.
\end{proof}

\begin{problem}\label{prob:terminal-cone}
For fixed $r$, determine the set
\[
 \mathcal N_r=\{n\ge r:\disc(J_{n,r})\text{ belongs to the square-graph cone}\}.
\]
The present results give
\[
 \mathcal N_2=\{n\ge2\},\qquad
 \mathcal N_3=\{n\ge3\},\qquad
 \mathcal N_4=\{4,5,\ldots,22\}.
\]
The first unknown case is $r=5$, for which
Proposition~\ref{prop:quintic-terminal} supplies an explicit starting point.
\end{problem}

\section{Finite graph types and the geometry of the cone}

The quartic threshold suggests separating a basic finiteness statement from
any positivity expectation.

\begin{proposition}[finite reduced graph types]\label{prop:finite-types}
Fix an even degree $d$.  There is a finite collection $\mathcal H_d$ of
square multigraphs without isolated vertices, each having at most $d$
vertices, such that for every $n\ge d$ the square-graph cone in
$\mathrm{PST}_{n,d}$ is generated by the symmetrizations of the graphs in
$\mathcal H_d$ after adjoining isolated vertices.  In the terminal problem of
order $r$ one may take $d=r(r-1)$.
\end{proposition}

\begin{proof}
If a square multigraph has total edge multiplicity $d$, every occupied
unordered pair of vertices has multiplicity at least two.  Hence there are at
most $d/2$ occupied pairs and at most $d$ nonisolated vertices.  After deleting
isolated vertices, only finitely many multigraphs can therefore occur.  For
$n\ge d$ every such reduced graph can be enlarged to $n$ vertices by adjoining
isolated vertices, which proves the claim.
\end{proof}

This elementary observation is compatible with the stable square-graph bases
constructed in \cite{AS}.  Theorem~\ref{thm:quartic-threshold} shows, however,
that finite graph type does \emph{not} force eventual positive membership.  In
fixed degree, cone membership becomes a finite-dimensional parametric linear
problem, and its dual can be equally effective: the single family
$\Lambda_n$ separates the terminal quartic discriminant for every $n\ge23$.

This motivates two questions which replace the uniform positivity expectation.

\begin{problem}[eventual cone behavior]\label{prob:eventual-cone}
For fixed $r$, determine whether membership of $\disc(J_{n,r})$ in the
square-graph cone is eventually constant as a function of $n$.  If it is,
determine the eventual sign and an effective threshold.
\end{problem}

\begin{problem}[stable dual certificates]\label{prob:dual-certificates}
For fixed degree $d$, describe rational or polynomial families of linear
functionals on $\mathrm{PST}_{n,d}$ which are nonnegative on every reduced
square-graph ray for all sufficiently large $n$.  In particular, find a
conceptual explanation of the quartic functional
\eqref{eq:quartic-separator} and of the face on which its $53$ zero graph types
lie.
\end{problem}

Equation~\eqref{eq:orbit-monomial-coordinate} makes these questions amenable
to exact finite computation: after a fixed shift in $n$, positivity of a dual
functional can sometimes be certified simply by coefficientwise positivity of
finitely many univariate polynomials.  This is precisely what happens for the
quartic obstruction.

\section{Relation with the square-graph cone}

Theorem~\ref{thm:quartic-threshold} shows that the historical conjecture fails
at terminal order four.  By contrast, the quadratic and cubic terminal
families remain uniformly inside the cone.  The cubic case is subtler than the
quadratic one: the formula of Theorem~\ref{thm:cubic} is not itself a
square-graph decomposition, because it contains the difference of two
invariant terms.  The essential point is therefore to prove that
\[
        p_2^3-\frac{n(n-1)}{(n-2)^2}p_3^2
\]
belongs to the square-graph cone.  The next section gives an explicit positive
certificate.  This certificate is stronger than the corresponding moment
inequality, because every summand is a symmetrized monomial in squared root
differences.


\section{A positive square-graph expansion in the cubic terminal case}

We now give the announced certificate.  This proves
Conjecture~\ref{conj:squaregraph} for the terminal cubic family
$k=n-3$.

Continue to use the centered variables
\[
 z_i=x_i-\bar x,
 \qquad \sum_i z_i=0,
 \qquad p_r=\sum_i z_i^r .
\]
Introduce the following four symmetric sums of square-graph monomials:
\[
\begin{aligned}
T_n&=\sum_{1\le i<j<k\le n}
 (z_i-z_j)^2(z_j-z_k)^2(z_k-z_i)^2,\\
P_n&=\sum_{i,j,k,l}^{*}
 (z_i-z_j)^2(z_j-z_k)^2(z_k-z_l)^2,\\
Q_n&=\sum_{i,j,k,l}^{*}
 (z_i-z_j)^4(z_k-z_l)^2,\\
R_n&=\sum_{i,j,k,l,m,s}^{*}
 (z_i-z_j)^2(z_k-z_l)^2(z_m-z_s)^2.
\end{aligned}
\]
Here $\sum^*$ means summation over ordered pairwise distinct indices.  Thus
$T_n$ is the triangle square-graph sum, $P_n$ is the ordered path sum,
$Q_n$ is the ordered sum of a double edge disjoint from a single edge, and
$R_n$ is the ordered matching sum.  Each of these belongs to the square-graph
cone.

\begin{lemma}\label{lem:degree-six-sums}
For centered variables $z_1,\ldots,z_n$ one has
\[
\begin{aligned}
T_n={}&-p_2^3-np_3^2+np_2p_4,\\
P_n={}&(4n-6)p_2^3-2(n-5)(n-2)p_3^2\\
&{} +(4n^2-24n+30)p_2p_4-2n(n-1)p_6,\\
Q_n={}&12(n-2)p_2^3+8(2n-5)p_3^2\\
&{} +(4n^2-36n+60)p_2p_4-4n(n-1)p_6,\\
R_n={}&8(n-2)(n^2-7n+15)p_2^3\\
&{} -16(3n^2-15n+20)p_3^2\\
&{} -24(n-2)(n^2-5n+10)p_2p_4\\
&{} +16n(n-1)(n-2)p_6 .
\end{aligned}
\]
\end{lemma}

\begin{proof}
We record the diagonal-removal calculation.  It is a useful check on the
normalization of the ordered sums.  For nonnegative integers
$\alpha_1,\ldots,\alpha_m$, put
\[
S(\alpha_1,\ldots,\alpha_m)=
\sum_{i_1,\ldots,i_m}^{*} z_{i_1}^{\alpha_1}\cdots z_{i_m}^{\alpha_m}.
\]
By inclusion--exclusion, or equivalently by M\"obius inversion on the lattice
$\Pi_m$ of set partitions, one has
\begin{equation}\label{eq:mobius-distinct}
\begin{aligned}
S(\alpha_1,\ldots,\alpha_m)
&=\sum_{\pi\in\Pi_m}\mu(\pi)
  \prod_{B\in\pi} p_{\sum_{j\in B}\alpha_j},\\
\mu(\pi)&=\prod_{B\in\pi}(-1)^{|B|-1}(|B|-1)! .
\end{aligned}
\end{equation}
Here $p_0=n$ and $p_1=0$.  Formula~\eqref{eq:mobius-distinct} is just the
standard principle that the unrestricted sum is decomposed according to the
partition of positions on which the indices coincide.

Expand each of the four square-graph products as a sum of monomials in the
corresponding ordered variables and apply~\eqref{eq:mobius-distinct}.  For the
triangle one first computes the ordered sum and divides by $6$, because each
unordered triple is counted six times; this gives
\[
T_n=-p_2^3-np_3^2+np_2p_4.
\]
The same expansion for the ordered path, the ordered disjoint double-edge plus
single-edge graph, and the ordered matching gives the displayed expressions for
$P_n,Q_n$ and $R_n$.  No specialization of the variables is used in this
calculation.
\end{proof}

\begin{theorem}\label{thm:cubic-square-expansion}
For every $n\ge3$,
\[
\begin{aligned}
 p_2^3-\frac{n(n-1)}{(n-2)^2}p_3^2
={}&\frac{6}{n(n-2)}T_n
 +\frac{1}{2n(n-2)}P_n\\
&{} +\frac{1}{4n(n-2)}Q_n
 +\frac{1}{8n(n-2)^2}R_n .
\end{aligned}
\]
Consequently $\D_{n-3}^{(n)}$ belongs to the square-graph cone for every
$n\ge3$.
\end{theorem}

\begin{proof}
Substitute the four identities of Lemma~\ref{lem:degree-six-sums} into the
right-hand side.  The coefficients of $p_2p_4$ and $p_6$ cancel.  The
coefficient of $p_2^3$ is
\[
-\frac{6}{n(n-2)}+\frac{4n-6}{2n(n-2)}
+\frac{12(n-2)}{4n(n-2)}
+\frac{8(n-2)(n^2-7n+15)}{8n(n-2)^2}=1,
\]
and the coefficient of $p_3^2$ is
\[
-\frac{6}{n-2}-\frac{n-5}{n}
+\frac{2(2n-5)}{n(n-2)}
-\frac{2(3n^2-15n+20)}{n(n-2)^2}
=-\frac{n(n-1)}{(n-2)^2}.
\]
This proves the displayed identity.  All four coefficients in the identity are
positive for $n\ge3$, and $T_n,P_n,Q_n,R_n$ are square-graph sums.  Combining
this identity with Theorem~\ref{thm:cubic} proves the square-graph cone
conjecture for $k=n-3$.
\end{proof}


\medskip
The preceding uniform identity can be compressed when $n\ge5$.  Define the
five-vertex square-graph sum
\[
U_n=\sum_{i,j,k,l,m}^{*}
(z_i-z_j)^2(z_j-z_k)^2(z_l-z_m)^2,
\]
where all five indices are pairwise distinct.  Thus $U_n$ is the ordered sum
associated with a doubled two-edge path together with a disjoint doubled
edge.

\begin{lemma}\label{lem:five-vertex-sum}
For centered variables one has
\begin{align*}
U_n={}&2(n-4)\Big(3(n-2)p_2^3+2(2n-5)p_3^2\\
&\hspace{2.4cm} +(n^2-9n+15)p_2p_4-n(n-1)p_6\Big).
\end{align*}
Equivalently,
\begin{equation}\label{eq:U-relation}
U_n=(n-4)\left(\frac12P_n+\frac14Q_n-(n-3)T_n\right).
\end{equation}
\end{lemma}

\begin{proof}
Apply the distinct-index M\"obius formula
\eqref{eq:mobius-distinct} to
$(z_i-z_j)^2(z_j-z_k)^2(z_l-z_m)^2$.  This gives the first displayed
formula.  Substitution of the formulas for $T_n,P_n,Q_n$ from
Lemma~\ref{lem:degree-six-sums} then gives \eqref{eq:U-relation}.
\end{proof}

\begin{theorem}[three-graph cubic certificate]\label{thm:cubic-threegraph}
For every $n\ge5$,
\begin{equation}\label{eq:cubic-threegraph}
\boxed{
 p_2^3-\frac{n(n-1)}{(n-2)^2}p_3^2
=
\frac{n+3}{n(n-2)}T_n
+\frac{1}{n(n-2)(n-4)}U_n
+\frac{1}{8n(n-2)^2}R_n .}
\end{equation}
All three coefficients are positive.  In particular, the terminal cubic
certificate uses only three square-graph types for $n\ge5$.
\end{theorem}

\begin{proof}
In Theorem~\ref{thm:cubic-square-expansion}, the $P_n$ and $Q_n$ terms occur
in the combination
\[
\frac{1}{n(n-2)}\left(\frac12P_n+\frac14Q_n\right).
\]
Using \eqref{eq:U-relation}, this equals
\[
\frac{U_n}{n(n-2)(n-4)}+\frac{n-3}{n(n-2)}T_n.
\]
Adding the existing coefficient $6/(n(n-2))$ of $T_n$ gives
$(n+3)/(n(n-2))$, while the $R_n$ coefficient is unchanged.
\end{proof}

\begin{corollary}[the Alexandersson--Shapiro three-graph formula]
\label{cor:AS-threegraph}
Put $n=k+3$.  For $k\ge3$, let $g_1,g_2,g_3$ be respectively the doubled
triangle, the disjoint union of a doubled two-edge path and a doubled edge,
and the matching of three doubled edges, with isolated vertices adjoined up
to $k+3$ vertices.  Then
\begin{align*}
\D_k^{(k+3)}=(k!)^3\Bigg[&
\frac{(k+1)^3(k+2)(k+6)}{72}\,\widetilde g_1
+\frac{(k+1)^3k(k+2)}{12}\,\widetilde g_2\\
&+\frac{(k-1)k(k+1)^2(k+2)(k-2)}{96}\,\widetilde g_3
\Bigg].
\end{align*}
For $k=2$ the same formula holds with the last (zero) term omitted.  Thus the
three-graph formula conjectured in \cite[Example~3]{AS} is proved throughout
its literal graph-theoretic range.
\end{corollary}

\begin{proof}
For $n\ge6$, the symmetrization convention of Section~2 gives
\[
\widetilde g_1=6(n-3)!\,T_n,\qquad
\widetilde g_2=(n-5)!\,U_n,\qquad
\widetilde g_3=(n-6)!\,R_n.
\]
Combine Theorems~\ref{thm:cubic} and \ref{thm:cubic-threegraph}, substitute
$n=k+3$, and simplify the factorials.  For $k=2$ the matching coefficient
vanishes and the first two identities suffice.
\end{proof}

\begin{remark}[a graph-theoretic proof of the equality case]
The positive certificates also give a direct proof of
Proposition~\ref{prop:equality}.  Indeed, if the cubic moment inequality is an
equality, then every graph sum occurring with positive coefficient must
vanish.  Since $T_n$ is a sum of squared Vandermonde products over triples,
$T_n=0$ implies that the centered roots assume at most two distinct values.
For $n\ge5$, if both multiplicities were at least two, one of them would be at
least three; choosing an alternating doubled two-edge path using two vertices
from the larger class and one from the smaller, together with a disjoint edge
joining the two classes, gives a strictly positive summand of $U_n$.  Hence one
multiplicity is $1$.  For $n=4$ the same conclusion follows from the $P_n$ term
in Theorem~\ref{thm:cubic-square-expansion}, and for $n=3$ it is automatic.
Thus the square-graph certificate detects the sharp equality configurations
without an optimization argument.
\end{remark}

\begin{corollary}\label{cor:cubic-discriminant-squaregraph}
For every $n\ge3$,
\[
\begin{aligned}
\D_{n-3}^{(n)}
={}&\frac{n!((n-2)!)^3}{12}
\left(
\frac{6}{n(n-2)}T_n
+\frac{1}{2n(n-2)}P_n
\right.\\
&\left.\hspace{2.7cm}
+\frac{1}{4n(n-2)}Q_n
+\frac{1}{8n(n-2)^2}R_n
\right).
\end{aligned}
\]
In particular, the terminal cubic discriminant has a manifestly nonnegative
square-graph expansion.
\end{corollary}

\begin{corollary}[first terminal cases of the square-graph conjecture]
\label{cor:first-terminal-cases}
Conjecture~\ref{conj:squaregraph} holds for $\D_{n-2}^{(n)}$ for all
$n\ge2$ and for $\D_{n-3}^{(n)}$ for all $n\ge3$.
\end{corollary}

\begin{proof}
The quadratic case is Proposition~\ref{prop:quadratic}.  The cubic case is
Corollary~\ref{cor:cubic-discriminant-squaregraph}; all coefficients in its
expansion are positive for $n\ge3$, and all four displayed graph sums are
square-graph sums.
\end{proof}

\section{Further directions}

The quartic threshold changes the natural continuation of the problem.  Some
concrete directions are now the following.

\begin{enumerate}[label={\rm(\arabic*)}]
\item Give a conceptual, non-computational explanation of the separating
functional \eqref{eq:quartic-separator}.  In particular, determine whether it
comes from a distinguished representation-theoretic quotient, a moment
functional, or an extremal ray of the dual square-graph cone.  Understanding
the $53$ graph types on which it vanishes may identify the relevant face.

\item Determine $\mathcal N_5$.  Proposition~\ref{prop:quintic-terminal}
reduces this to an explicit degree-$20$ problem, while
Corollary~\ref{cor:quintic-odd-slice} completely factors the first natural
slice $p_3=p_5=0$.  The quartic result suggests that one should search
simultaneously for positive certificates at small $n$ and stable dual
separators at large $n$, rather than assuming one behavior in advance.

\item Decide whether eventual cone behavior in
Problem~\ref{prob:eventual-cone} holds for every fixed terminal order $r$.
Equation~\eqref{eq:orbit-monomial-coordinate} turns this into a finite family
of parametric cone problems, but the geometry can change with $n$.

\item Although $\Delta_{n,4}$ leaves the square-graph cone for $n\ge23$, it is
still nonnegative on real root configurations.  Determine whether it belongs
to a larger natural sum-of-squares cone in the differences, and compare this
with the sum-of-squares theorem for first derivatives in \cite{Sanyal}.

\item Determine whether the terminal Appell chain has a direct
representation-theoretic or probabilistic explanation that also controls the
dual cones.  The subset-average formula \eqref{eq:subset-average} suggests that
finite-population sampling operators may be relevant.
\end{enumerate}

\begin{thebibliography}{99}

\bibitem{AS}
\textsc{P.~Alexandersson and B.~Shapiro},
Discriminants, symmetrized graph monomials, and sums of squares,
\emph{Experiment. Math.} \textbf{21} (2012), no.~4, 353--361.

\bibitem{AScert}
\textsc{P.~Alexandersson},
Discriminants, symmetrized graph monomials, and sums of squares certificates,
arXiv:1202.0416.

\bibitem{GKZ}
\textsc{I.~M.~Gel'fand, M.~M.~Kapranov and A.~V.~Zelevinsky},
\emph{Discriminants, Resultants and Multidimensional Determinants},
Birkh\"auser, Boston, 1994.

\bibitem{Janson}
\textsc{S.~Janson},
Resultant and discriminant of polynomials,
\emph{unpublished notes}, 2007.

\bibitem{Sanyal}
\textsc{R.~Sanyal, B.~Sturmfels and C.~Vinzant},
The entropic discriminant,
\emph{Adv. Math.} \textbf{244} (2013), 678--707.

\bibitem{So1}
\textsc{F.~Sottile},
Frontiers of reality in Schubert calculus,
\emph{Bull. Amer. Math. Soc.} \textbf{47} (2010), 31--71.

\end{thebibliography}

\end{document}
